\section{Shifting natural-number indices}

\import{set/finite-ordinal.tex}
\import{function.tex}

Translation of natural indices also identifies a finite block with the tail
of a larger finite ordinal. The same translation graph is used for sequences.

\begin{definition}\label{natural_shift}
    $\naturalShift{N}=\{(n,N+n)\mid n\in\naturals\}$.
\end{definition}

\begin{theorem}\label{natural_shift_is_function}
    Let $N$ be a natural number.
    Then $\naturalShift{N}\in\funs{\naturals}{\naturals}$.
\end{theorem}
\begin{proof}
    $\naturalShift{N}$ is a relation by
        \auto{natural_shift,relation}.
    $\naturalShift{N}$ is right-unique by
        \auto{natural_shift,rightunique,pair_eq_iff}.
    $\dom{\naturalShift{N}}=\naturals$ by
        \auto{natural_shift,setext,dom_iff,pair_eq_iff}.
    For every $n\in\naturals$ we have
        $(\naturalShift{N})(n)=N+n$ by
        \obv{natural_shift,function_apply_intro}.
    For every $n\in\naturals$ we have $N+n\in\naturals$ by
        \obv{num_addition_on_naturals}.
    For every $n\in\dom{\naturalShift{N}}$ we have
        $(\naturalShift{N})(n)\in\naturals$ by \auto.
    Follows by \obv{funs_intro}.
\end{proof}

\begin{proposition}\label{natural_shift_apply}
    Let $N,n$ be natural numbers.
    Then $(\naturalShift{N})(n)=N+n$.
\end{proposition}
\begin{proof}
    $\naturalShift{N}$ is a function by
        \auto{natural_shift_is_function,funs_is_function}.
    Follows by \obv{natural_shift,function_apply_intro}.
\end{proof}

\subsection{Translation preserves order}

\begin{theorem}\label{natural_addition_subseteq_iff}
    Let $m,n,k$ be natural numbers.
    Then $m+n\subseteq m+k$ iff $n\subseteq k$.
\end{theorem}
\begin{proof}
    $m+n$ and $m+k$ are natural numbers by \obv{num_addition_on_naturals}.
    Show that if $m+n\subseteq m+k$, then $n\subseteq k$.
    \begin{subproof}
        Assume $m+n\subseteq m+k$.
        Take a natural number $d$ such that $m+k=(m+n)+d$ by
            \obv{natural_additive_difference}.
        $(m+n)+d=m+(n+d)$ by \obv{natural_addition_associative}.
        $n+d$ is a natural number by \obv{num_addition_on_naturals}.
        Hence $k=n+d$ by \auto{natural_addition_cancel_left}.
        Follows by \auto{natural_addition_left_subseteq_sum}.
    \end{subproof}
    Show that if $n\subseteq k$, then $m+n\subseteq m+k$.
    \begin{subproof}
        Assume $n\subseteq k$.
        Take a natural number $d$ such that $k=n+d$ by \obv{natural_additive_difference}.
        $m+k=(m+n)+d$ by \auto{natural_addition_associative}.
        Follows by \auto{natural_addition_left_subseteq_sum}.
    \end{subproof}
    Follows by \auto.
\end{proof}

\begin{theorem}\label{natural_addition_membership_iff}
    Let $m,n,k$ be natural numbers.
    Then $m+n\in m+k$ iff $n\in k$.
\end{theorem}
\begin{proof}
    $m+n$ and $m+k$ are natural numbers by \obv{num_addition_on_naturals}.
    $m+n$ and $m+k$ are ordinals by \obv{natural_number_is_ordinal}.
    $n$ and $k$ are ordinals by \obv{natural_number_is_ordinal}.
    $m+n\subseteq m+k$ iff $n\subseteq k$ by \obv{natural_addition_subseteq_iff}.
    $m+n=m+k$ iff $n=k$ by \auto{natural_addition_cancel_left}.
    Follows by \auto{ordinal_proper_subset_implies_elem,ordinal_elem_implies_proper_subset,subset}.
\end{proof}

\subsection{Finite blocks}

\begin{definition}\label{finite_index_shift}
    $\finiteIndexShift{m}{n}=\restrl{\naturalShift{m}}{n}$.
\end{definition}

\begin{proposition}\label{finite_index_shift_apply}
    Let $m,n$ be natural numbers and $k\in n$.
    Then $(\finiteIndexShift{m}{n})(k)=m+k$.
\end{proposition}
\begin{proof}
    $n\subseteq\naturals$ by \auto{naturals_is_transitive_set,transitiveset_iff_subseteq}.
    Hence $k\in\naturals$ by \obv{subseteq}.
    $\naturalShift{m}\in\funs{\naturals}{\naturals}$ by \obv{natural_shift_is_function}.
    Follows by \auto{finite_index_shift,restrl_of_function_apply,natural_shift_apply,funs_is_function,funs}.
\end{proof}

\begin{theorem}\label{finite_index_shift_type}
    Let $m,n$ be natural numbers.
    Then $\finiteIndexShift{m}{n}\in\funs{n}{(m+n)\setminus m}$.
\end{theorem}
\begin{proof}
    Let $e=\finiteIndexShift{m}{n}$.
    $n\subseteq\naturals$ by \auto{naturals_is_transitive_set,transitiveset_iff_subseteq}.
    $e\in\funs{n}{\naturals}$ by \auto{finite_index_shift,natural_shift_is_function,funs_restrl}.
    Hence $e$ is a function and $\dom{e}=n$ by \obv{funs_is_function,funs}.
    Show that for every $k\in n$ we have $e(k)\in(m+n)\setminus m$.
    \begin{subproof}
        Fix $k\in n$.
        $e(k)=m+k$ by \auto{finite_index_shift_apply}.
        $m+k\in m+n$ by \obv{natural_addition_membership_iff}.
        $m\subseteq m+k$ by \obv{natural_addition_left_subseteq_sum}.
        Hence $m+k\notin m$ by \auto{subseteq,in_irrefl}.
        Follows by \obv{setminus}.
    \end{subproof}
    Hence $e$ is a function to $(m+n)\setminus m$.
    Follows by \auto{funs_intro}.
\end{proof}

\begin{theorem}\label{finite_index_shift_bijection}
    Let $m,n$ be natural numbers.
    Then $\finiteIndexShift{m}{n}\in\bijections{n}{(m+n)\setminus m}$.
\end{theorem}
\begin{proof}
    Let $e=\finiteIndexShift{m}{n}$.
    $n\subseteq\naturals$ by \auto{naturals_is_transitive_set,transitiveset_iff_subseteq}.
    $e\in\funs{n}{(m+n)\setminus m}$ by \auto{finite_index_shift_type}.
    Hence $e$ is a function and $\dom{e}=n$ by \obv{funs_is_function,funs}.
    For all $k,l\in n$ such that $e(k)=e(l)$ we have $k=l$ by
        \auto{finite_index_shift_apply,natural_addition_cancel_left,subseteq}.
    Hence $e$ is injective by \auto{injective_function}.
    Show that for every $v\in(m+n)\setminus m$ there exists $k\in n$ such that $e(k)=v$.
    \begin{subproof}
        Fix $v\in(m+n)\setminus m$.
        $v\in m+n$ and $v\notin m$ by \obv{setminus}.
        $m+n$ is a natural number by \obv{num_addition_on_naturals}.
        $m+n\subseteq\naturals$ by \auto{naturals_is_transitive_set,transitiveset_iff_subseteq}.
        Hence $v$ is a natural number by \obv{subseteq}.
        $m$ and $v$ are ordinals by \obv{natural_number_is_ordinal}.
        Hence $m\subseteq v$ by \obv{ordinal_elem_or_subseteq}.
        Take a natural number $k$ such that $v=m+k$ by \obv{natural_additive_difference}.
        Hence $k\in n$ by \auto{natural_addition_membership_iff}.
        $e(k)=v$ by \auto{finite_index_shift_apply}.
    \end{subproof}
    Hence $e\in\surj{n}{(m+n)\setminus m}$ by \obv{surj}.
    Follows by \auto{bijections}.
\end{proof}

\begin{proposition}\label{finite_index_shift_zero}
    Let $n$ be a natural number.
    Then $\finiteIndexShift{\emptyset}{n}=\identity{n}$.
\end{proposition}
\begin{proof}
    $\emptyset$ is a natural number by \obv{omega_empty}.
    $\emptyset+n=n$ by \obv{num_addition_zero_left}.
    $\finiteIndexShift{\emptyset}{n}\in\funs{n}{n}$ by
        \auto{finite_index_shift_type,setminus_emptyset}.
    $\identity{n}\in\funs{n}{n}$ by \obv{id_elem_funs}.
    $n\subseteq\naturals$ by \auto{naturals_is_transitive_set,transitiveset_iff_subseteq}.
    For every $k\in n$ we have $(\finiteIndexShift{\emptyset}{n})(k)=(\identity{n})(k)$ by
        \auto{finite_index_shift_apply,num_addition_zero_left,id_apply,subseteq}.
    Follows by \obv{typed_functions_equal_from_agreement}.
\end{proof}

\begin{proposition}\label{natural_sum_member_iff}
    Let $m,n$ be natural numbers.
    Then $i\in m+n$ iff $i\in m$ or there exists $k\in n$ such that $i=m+k$.
\end{proposition}
\begin{proof}
    Let $e=\finiteIndexShift{m}{n}$ and $D=(m+n)\setminus m$.
    $e\in\bijections{n}{D}$ by \auto{finite_index_shift_bijection}.
    Hence $e\in\funs{n}{D}$ and $e\in\surj{n}{D}$ by \obv{bijections_to_funs,bijections}.
    $m\subseteq m+n$ by \obv{natural_addition_left_subseteq_sum}.
    Show that if $i\in m+n$, then $i\in m$ or there exists $k\in n$ such that $i=m+k$.
    \begin{subproof}
        Assume $i\in m+n$.
        \begin{byCase}
            \caseOf{$i\in m$.}
                Follows.
            \caseOf{$i\notin m$.}
                $i\in D$ by \auto{setminus}.
                Take $k\in n$ such that $e(k)=i$ by \obv{surj}.
                $i=m+k$ by \auto{finite_index_shift_apply}.
        \end{byCase}
    \end{subproof}
    Show that if $i\in m$ or there exists $k\in n$ such that $i=m+k$, then $i\in m+n$.
    \begin{subproof}
        Assume $i\in m$ or there exists $k\in n$ such that $i=m+k$.
        \begin{byCase}
            \caseOf{$i\in m$.}
                Follows by \obv{subseteq}.
            \caseOf{There exists $k\in n$ such that $i=m+k$.}
                Take $k\in n$ such that $i=m+k$.
                $e(k)=i$ by \auto{finite_index_shift_apply}.
                $e(k)\in D$ by \obv{funs_type_apply}.
                Follows by \auto{setminus}.
        \end{byCase}
    \end{subproof}
    Follows by \auto.
\end{proof}
